\chapter{Limiti di funzioni}
\label{limiti-di-funzioni}

\section{Intorni}

\defn{Intorni}{
    Sia \(x_{0}\in \mathbb{R}\), un Intorno \(I(x_{0})\) di \(x_{0}\) è un intervallo aperto contenente \(x_{0}\).
    \begin{itemize}
        \item Un intorno \textbf{Simmetrico} di \(x_{0}\) di ampiezza \(\delta>0\) è un intervallo del tipo:
        \[]x_{0}-\delta, x_{0} +\delta[ = \boxed{I_{\delta}(x_{0})}\]
        \item Un intorno \textbf{Sinistro} di \(x_{0}\) è un intervallo del tipo:
        \[]x_{1}, x_{0}[ \text{ con }x_{1} <x_{0} = \boxed{I^{-}(x_{0})}\]
        \item Un intorno \textbf{Destro} di \(x_{0}\) è un intervallo del tipo:
        \[]x_{0}, x_{2}[ \text{ con }x_{2} > x_{0} = \boxed{I^{+}(x_{0})}\]
    \end{itemize}
}

\section{Punti d'accumulazione}

\defn{Punti d'Accumulazione}{
    Se \(A \subseteq \mathbb{R}\) e sia \(\underline{x_{0}\in \mathbb{R}}\).
    Diremo che \(x_{0}\) è un \textbf{punto di accumulazione} per \(A\) se \(\forall I(x_{0}) \exists\) almeno un punto di \(A\) diverso da \(x_{0}\) ovvero:
    \[\forall I (x_{0}) \implies I(x_{0})\cap A \setminus \{ x_{0} \}\ne \varnothing \]
    Un modo equivalente è:
    \[\forall \epsilon>0 \exists \overline{x} \in A : 0 <|\overline{x} - x_{0}| < \epsilon \implies x_{0} - \epsilon < \overline{x} <x_{0} +\epsilon\]
}

\rmkb{
    \begin{itemize}
        \item \(+\infty\) è un pto di accumulazione per \(A \iff A\) non è superiormente limitato
        \item \(-\infty\) è un pto di accumulazione per \(A \iff A\) non è inferiormente limitato
    \end{itemize}
}

\ex{
    \begin{itemize}
        \item \(A = \left\{ \frac{1}{n}, n\in \mathbb{N} \right\} \quad x_{0}=0 \not\in A\)
        \item \(A=\mathbb{N} \quad +\infty\)
        \item \(A = ]0,1[ \qquad [0,1]\)
        \item \(A = \mathbb{Q} \qquad \overline{\mathbb{R}}\) (Per il lemma di densità di \(\mathbb{Q}\) in \(\mathbb{R}\))
    \end{itemize}
}

\noindent\rule{0.5\linewidth}{0.5pt}

\prop{
    Sia \(A \subseteq \mathbb{R}\) e \(x_{0}\in \mathbb{R}\).
    \[
    \underset{(1)}{\boxed{x_{0} \text{ è di acc. per A}}}\iff 
    \underset{(2)}{\boxed{\begin{array}{l}
    \exists x_{n} \text{ una successione}:\\
    x_{n} \in A \\
    x_{n} \neq x_{0} \quad \forall n \in \mathbb{N} \\
    x_{n} \to x_{0}
    \end{array}}}
    \]
}

\pf{
    \((\impliedby)(2)\)
    Per Hp:
    \[\forall \epsilon >0 \exists \nu: |x_{n}-x_{0}| < \epsilon \quad \forall n > \nu \implies x_{0} \text{ di acc}\]
    oppure
    \[\forall I(x_{0}) \exists \nu: x_{n} \in I(x_0) \quad \forall n > \nu\]
    
    \((\implies)(1)\)
    Per Hp:
    \[\forall \epsilon>0 \exists \overline{x} \in A : 0 < | \overline{x} - x_{0}| <\epsilon\]
    Preso \(\epsilon = \frac{1}{n}, n \in \mathbb{N}\):
    \[\boxed{\exists x_{n} \in A \forall n} \text{ e } \boxed{x_{n}\ne x_{0} \forall n:} 0< |x_{n} -x_{0} | < \frac{1}{n} \implies x_{n} \to x_{0}\]
}

\section{Teorema di Bolzano}
\label{teo-di-bolzano}

\thm{Teorema di Bolzano}{
    Ogni insieme \textbf{Limitato} e \textbf{infinito} di \(\mathbb{R}\) ammette almeno un punto di accumulazione.
}

\pf{
    Poiché l'insieme \(A\) è infinito, contiene una successione \(x_{n}\) di punti tutti diversi tra loro.
    Per il \textbf{Teorema di Bolzano-Weierstrass} (vedi capitolo Successioni) \(\exists x_{n_{k}}\in A: x_{n_{k}}\to l \in \mathbb{R}\) e \(l\) è di accumulazione.
}

\section{Insiemi Aperti e Chiusi e Compatto}

\defn{Aperto}{
    Un insieme \(A \subseteq \mathbb{R}\) si dice aperto se:
    \[\forall x_{0} \in A \exists I(x_{0})\subseteq A\]
}

\ex{
    \(]0,1[\) SÌ \\
    \([0,1[\) NO
}

\defn{Chiuso}{
    \(C \subseteq \mathbb{R}\) si dice chiuso se il suo complementare \(\mathbb{R}\setminus C\) è aperto.
}

\ex{
    \(C: [0,1]\) \\
    \(R\setminus C = ]-\infty,0[ \cup ]1, +\infty[\)
}

\prop{
    \[C \subseteq \mathbb{R} \text{ è chiuso} \iff \text{contiene tutti i suoi punti di accumulazione}\]
}

\pf{
    \textbf{\((\implies)\)}
    Se \(C\) è chiuso \(\implies R\setminus C=A\) è aperto.
    Sia \(x_{0}\in \mathbb{R}, x_{0}\) di accumulazione per \(C\).
    Se per assurdo \(x_{0}\not \in C\) allora \(x_{0} \in A\); essendo \(A\) aperto:
    \[\exists I(x_{0})\subseteq A \implies I(x_{0})\cap C = \varnothing \text{ ASSURDO}\]
    
    \textbf{\((\impliedby)\)}
    Sia \(A=\mathbb{R}\setminus C\). Se per assurdo \(A\) non è aperto, \(\exists x_{0}\in A:\forall n\in \mathbb{N}\), preso \(\epsilon=\frac{1}{n}\):
    \[\exists x_{n}:|x_{n}-x_{0}| <\frac{1}{n} \text{ ma }x_{n}\in C \; (x_{n}\ne x_{0}) \quad x_{n}\to x_{0}\in C\]
    ASSURDO. Allora \(A\) è aperto.
}

\prop{
    \[C \subseteq \mathbb{R} \text{ è chiuso} \iff \text{contiene i limiti di tutte le sue successioni}\]
}

\pf{
    Per esercizio.
}

\defn{Compatto}{
    \(K \subseteq \mathbb{R}\) si dice compatto se da ogni sua successione se ne può estrarre una convergente ad un elemento di \(K\).
}

\section{Teorema di Heine-Borel}

\thm{Teorema di Heine-Borel}{
    \(K \subseteq \mathbb{R}\) è compatto \(\iff K\) è chiuso e limitato.
}

\pf{
    \textbf{\((\impliedby)\)}
    Sia \(x_{n} \in K\).
    La limitatezza ci dice che \(x_{n_{k}} \to x_{0} \in \mathbb{R}\).
    La chiusura ci dice che \(\exists x_{n_{k}}\to x_{0} \in K\).

    \textbf{\((\implies)\)}
    Sia \(x_{n} \in K\) convergente a \(x_{0} \in K\). \(K\) è chiuso.
    Se per assurdo \(K\) non è limitato, allora \(\exists x_{n} \to +\infty\). Ma \(x_{n}\in K\) deve ammettere una sottosuccessione convergente ad un elemento di \(K\) (per l'ipotesi di compattezza), il che è assurdo.
}

\ex{
    \([1,2] \cup [3,5]\cup[6,7]\)
}

\section{Dai Limiti di Successioni a quelli di Funzioni}

\defn{Limite di Funzione}{
    Sia \(f:A \subseteq \mathbb{R} \to \mathbb{R}\) e sia \textbf{\(x_{0}\in \mathbb{R}\)} un punto di accumulazione per \(A\).
    Diremo che \(f\) converge o tende ad \textbf{\(l\in \mathbb{R}\)} per \(x \to x_{0}\) e scriveremo:
    \[\lim_{ x \to x_{0} } f(x)=l\]
    Se \(\forall x_{n} \in A\; \forall n, x_{n}\neq x_{0} \forall n, x_{n}\to x_{0}\) risulta:
    \[\lim_{ n \to +\infty } f(x_{n})=l\]
}

\ex{
    \(A= ]0, +\infty[\), \(x_{0}=0\) p.a.
    \[\forall x_{n}\in]0, +\infty[\quad x_{n}\to 0 \implies \lim_{ n \to +\infty } \log(x_{n})=-\infty\]
    % [Immagine: def_log_a-mag-1.jpg]
    % \begin{figure}[h] \centering \includegraphics[width=0.5\textwidth]{img/def_log_a-mag-1.jpg} \caption{Logaritmo} \end{figure}
}

\ex{
    \(f(x)=\frac{\sin x}{x} \quad A=\{ x\neq 0 \}= \mathbb{R} \setminus \{ 0 \} \quad x_{0}=0 \text{ è p.a.}\)
    \[
    \lim_{ x \to 0 } \frac{\sin x}{x}
    \]
    Sia \(x_{n} \in A, x_{n} \to 0\). Allora:
    \[
    \lim_{ x \to 0 } \frac{\sin x}{x}=\boxed{\lim_{ n \to +\infty } \frac{\sin x_{n}}{x_{n}} = 1}
    \]
}

\ex{
    Sia \(f: \mathbb{R} \to \mathbb{R}\) data da
    \[f(x)=\begin{cases}
    1 & \text{se } x\neq 0 \\
    0 & \text{se } x = 0
    \end{cases}\]
    Sia \(x_{n}\) una successione: \(x_{n} \in \mathbb{R}, x_{n}\neq 0 \forall n, x_{n}\to 0\).
    \[\lim_{ x \to 0 } f(x)=\lim_{ n \to +\infty } f(x_{n}) = 1 \neq f(0)\]
    % [Immagine: 1 tranne in 0.jpg]
}

\ex{
    \[f(x)=\begin{cases}
    0 & \text{se } x<0 \\
    1 & \text{se } x\geq 0
    \end{cases}\]
    \[\lim_{ x \to 0 } f(x) = \not \exists\]
    Se esistesse, \(\forall x_{n} \to 0, x_{n} \in \mathbb{R}, x_{n} \neq 0 \quad \forall n \implies \lim_{ n \to +\infty } f(x_{n})=l\).
    Ma se prendiamo:
    \[x_{n}=\frac{1}{n} \quad x'_{n}=-\frac{1}{n}\]
    \[\lim_{ n \to +\infty } f(x_{n}) =1 \neq \lim_{ n \to +\infty } f(x'_{n}) = 0\]
    Quindi \(\lim_{ x \to 0 } f(x) = \not \exists\).
    % [Immagine: funz gradino.jpg]
}

\rmkb{
    La stessa definizione vale nel caso in cui \(l \in \overline{\mathbb{R}}\) e \(x_{0}\in \overline{\mathbb{R}}\).
    \ex{
        \[f(x)=\frac{1}{x^2} \quad A= \mathbb{R} \setminus \{0\}\]
        Se \(x_{n}\in A\) e \(x_{n} \to \infty\):
        \[\lim_{ x \to +\infty } f(x) = \lim_{ n \to +\infty } \frac{1}{x_{n}^2}=0\]
        Se \(x_{n}\to -\infty\):
        \[\lim_{ x \to -\infty } f(x) = \lim_{ n \to +\infty } \frac{1}{x_{n}^2}=0\]
        Se \(x_{n}\in A, x_{n} \to 0, x_{n} \neq 0\):
        \[\lim_{ x \to 0 } f(x) = \lim_{ n \to +\infty } \frac{1}{x_{n}^2}=+\infty\]
        Se \(f(x)=\frac{1}{x} \implies \lim_{ x \to 0 } f(x)= \not \exists\)
    }
}

\defn{Limiti Infiniti}{
    Se \(f: A \subseteq \mathbb{R} \to \mathbb{R}, x_{0}\in \mathbb{R}\) di accumulazione.
    \begin{align*}
    & \lim_{ x \to x_{0} } f(x) = +\infty \iff \forall x_{n} \in A, x_{n}\to x_{0}, x_{n}\neq x_{0} \implies \lim_{ n \to +\infty } f(x_{n}) = +\infty\\
    & \lim_{ x \to x_{0} } f(x) = -\infty \iff \forall x_{n} \in A, x_{n}\to x_{0}, x_{n}\neq x_{0} \implies \lim_{ n \to +\infty } f(x_{n}) = -\infty\\
    \end{align*}
    \begin{itemize}
        \item Sia \(x_{0}=+\infty\) di acc per \(A\) (Non superiormente limitato) allora:
        \[\lim_{ x \to +\infty } f(x) = l \in \overline{\mathbb{R}} \iff \forall x_{n}\in A , x_{n} \to +\infty \implies \lim_{ n \to +\infty } f(x_{n})=l\]
        \item Sia \(x_{0}=-\infty\) di acc per \(A\) (Non inferiormente limitato) allora:
        \[\lim_{ x \to -\infty } f(x) = l \in \overline{\mathbb{R}} \iff \forall x_{n}\in A , x_{n} \to -\infty \implies \lim_{ n \to +\infty } f(x_{n})=l\]
    \end{itemize}
}

\noindent\rule{0.5\linewidth}{0.5pt}

\defn{Limite (Definizione Topologica)}{
    Sia \(f:A\subseteq \mathbb{R}\to \mathbb{R}\) e \(x_{0} \in \mathbb{R}\) di acc. per \(A\). Diremo che
    \[\lim_{ x \to x_{0} } f(x) =l \in \mathbb{R} \iff \forall \epsilon > 0 \exists \delta>0 : |f(x)-l|<\epsilon \; \forall x\in A , 0<|x-x_{0}|<\delta\]
    Oppure
    \[\iff \forall I(l)\exists I (x_{0}):f(x)\in I(l)\; \forall x \in A \cap I (x_{0})\setminus \{ x_{0} \}\]
}

\subsection{Teorema Ponte}

\thm{Teorema Ponte}{
    Sia \(f: A \subseteq \mathbb{R} \to \mathbb{R}\) e \(x_{0}\in \mathbb{R}\) di acc. per \(A\).
    Allora sono equivalenti le seguenti:
    \begin{align*}
    &(1) \quad \forall x_{n}\in A, x_{n}\neq x_{0} \forall n, x_{n}\to x_{0} \implies \lim_{ n \to +\infty } f(x_{n}) =l =\lim_{ x \to x_{0} } f(x)\\
    & \Updownarrow \\
    &(2) \quad \forall \epsilon > 0 \;\exists\delta>0 : |f(x)-l| < \epsilon \; \forall x \in A : 0< |x-x_{0}| < \delta \quad (\text{dove } l=\lim_{ x \to x_{0} } f(x))
    \end{align*}
    Dove (1)=Ipotesi e (2)=Tesi (e viceversa).
}

\pf{
    \textbf{(1)\(\implies\)(2)}
    Se per assurdo (2) non valesse:
    \[\exists \epsilon_{0}>0 : \forall \epsilon>0 \exists x_{\delta}: |f(x_{\delta})-l| \geq \epsilon_{0}\]
    dove \(0<|x_{\delta}-x_{0}|<\delta\).
    Se allora scelgo \(\delta=\frac{1}{n}\),
    \[\forall n \exists x_{n}: | f(x_{n})-l|\geq\epsilon_{0} \implies \boxed{0<|x_{n}-x_{0}|<\frac{1}{n}}\]
    Ma questo implica \(x_n \to x_0\) e \(f(x_n)\) non tende a \(l\), contraddicendo (1).

    \textbf{(2)\(\implies\)(1)}
    Sia \(x_{n}\to x_{0},x_{n}\neq x_{0}, x_{n}\in A\).
    Dobbiamo dimostrare che \(\lim_{n \to +\infty} f(x_n) = l\), cioè:
    \[\forall \epsilon > 0 \exists \overline{\nu} : |f(x_{n})-l| < \epsilon \; \forall n > \overline{\nu}\]
    Per ipotesi (2): \(\forall \epsilon > 0 \exists \delta > 0 : |f(x)-l| < \epsilon\) quando \(0 < |x - x_0| < \delta\).
    Poiché \(x_n \to x_0\), \(\exists \nu\) tale che per \(n > \nu\) si ha \(0 < |x_n - x_0| < \delta\) (essendo \(x_n \ne x_0\)).
    Quindi per tali \(n\) vale \(|f(x_n) - l| < \epsilon\).
}

\subsection{Le Nove Definizioni Di Limite}

\begin{enumerate}
    \item \(\lim_{ x \to x_{0} } f(x)=l \iff \forall \epsilon > 0 \exists \delta > 0 : |f(x)-l|<\epsilon \; \forall x \in A : 0 < |x -x_{0}| < \delta\)
    \item \(\lim_{ x \to x_{0} } f(x)=+\infty \iff \forall M > 0 \exists \delta > 0 : f(x)>M\; \forall x \in A : 0 < |x -x_{0}| < \delta\)
    \item \(\lim_{ x \to x_{0} } f(x)=-\infty \iff \forall M > 0 \exists \delta > 0 : f(x)<-M\; \forall x \in A : 0 < |x -x_{0}| < \delta\)
    \item \(\lim_{ x \to +\infty } f(x)=l \iff \forall \epsilon>0 \exists M>0 : |f(x)-l|<\epsilon\; \forall x \in A , x>M\)
    \item \(\lim_{ x \to +\infty } f(x)=+\infty \iff \forall K>0 \exists M>0 : f(x)>K \; \forall x \in A , x>M\)
    \item \(\lim_{ x \to +\infty } f(x)=-\infty \iff \forall K>0 \exists M>0 : f(x)<-K \; \forall x \in A , x>M\)
    \item \(\lim_{ x \to -\infty } f(x)=l \iff \forall \epsilon>0 \exists M>0 : |f(x)-l|<\epsilon\; \forall x \in A , x<-M\)
    \item \(\lim_{ x \to -\infty } f(x)=+\infty \iff \forall K>0 \exists M>0 : f(x)>K \; \forall x \in A , x<-M\)
    \item \(\lim_{ x \to -\infty } f(x)=-\infty \iff \forall K>0 \exists M>0 : f(x)<-K \; \forall x \in A , x<-M\)
\end{enumerate}

\subsection{Teorema Di Unicità Del Limite}

\prop{
    Se \(\lim_{ x \to x_{0} }f(x) = l \in \overline{\mathbb{R}}\)
    dove \(f:A \subseteq \mathbb{R} \to \mathbb{R}\) e \(x_{0}\in \overline{\mathbb{R}}\) è di acc. per A, Allora \(l\) è unico.
}

\pf{
    Segue dal \textbf{Teorema Ponte} e dal Teorema di Unicità del limite per le successioni:
    \[\forall x_{n} \in A: x_{n}\to x_{0}, x_{n}\neq x_{0} \implies f(x_{n})\to l\]
    Poiché il limite di successione è unico, anche il limite di funzione è unico.
}

\subsection{Operazioni Con I Limiti}

\prop{
    Siano \(f,g: A \subseteq \mathbb{R} \to \mathbb{R}\) e \(x_{0} \in \overline{\mathbb{R}}\) di acc per \(A\).
    Sia \(\lim_{ x \to x_{0} } f(x)=l_{1}\) e \(\lim_{ x \to x_{0} } g(x)=l_{2}\) con \(l_{1},l_{2} \in \overline{\mathbb{R}}\).
    Allora (salvo forme indeterminate):
    \begin{enumerate}
        \item \(\lim_{ x \to x_{0} } (f(x)\pm g(x))=l_{1}\pm l_{2}\)
        \item \(\lim_{ x \to x_{0} } (f(x)g(x))=l_{1}l_{2}\)
        \item \(\lim_{ x \to x_{0} } \frac{f(x)}{g(x)}=\frac{l_{1}}{l_{2}} \quad (g(x)\neq 0 \; \forall x \in A, l_2 \neq 0)\)
        \item \(\lim_{ x \to x_{0} } [f(x)]^{g(x)} \quad (f(x)>0)\)
    \end{enumerate}
}

\pf{
    Segue dal \textbf{Teorema Ponte} applicando le proprietà corrispondenti per le successioni.
}

\subsection{Teorema Del Confronto}

\defn{Teorema del Confronto}{
    Siano \(f(x),g(x),h(x)\) tre funzioni definite in \(A\) e sia \(x_{0} \in \overline{\mathbb{R}}\) di acc. per \(A\).
    \begin{enumerate}
        \item Se \(f(x)\leq g(x)\leq h(x)\) \(\forall x\) in un intorno di \(x_{0}\) e se:
        \[\lim_{ x \to x_{0} }f(x) =l \in \mathbb{R} \quad \text{e} \quad \lim_{ x \to x_{0} }h(x) =l \in \mathbb{R}\]
        Allora:
        \[\lim_{ x \to x_{0} } g(x)=l\]
        \item Se \(f(x)\leq g(x)\) in un intorno di \(x_{0}\):
        \[\lim_{ x \to x_{0} } f(x)=+\infty \implies \lim_{ x \to x_{0} } g(x)=+\infty\]
        \item Se \(g(x)\leq f(x)\) in un intorno di \(x_{0}\):
        \[\lim_{ x \to x_{0} } f(x)=-\infty \implies \lim_{ x \to x_{0} } g(x)=-\infty\]
    \end{enumerate}
}

\subsection{Teorema Del Limite Delle Funzioni Composte}

\prop{
    Siano \(f,g\) tali che sia definita \(f \circ g\).
    Siano \(x_{0},y_{0}\) punti di accumulazione dei domini di \(f,g\) rispettivamente.
    Supponiamo inoltre:
    \begin{enumerate}
        \item \(\lim_{ y \to y_{0} } g(y)=x_{0}\)
        \item \(\lim_{ x \to x_{0} } f(x)=l\in \mathbb{R}\)
        \item (Ipotesi aggiuntiva necessaria per il teorema standard) \(g(y) \neq x_0\) in un intorno di \(y_0\) (eccetto al più \(y_0\)) OPPURE \(f\) continua in \(x_0\).
    \end{enumerate}
    Allora:
    \[\lim_{ y \to y_{0} } f(g(y))=l\]
}

\pf{
    Per il \textbf{Teorema Ponte}.
    Sia \(y_{n}\to y_{0}\) con \(y_{n}\neq y_{0} \; \forall n\).
    Poniamo \(x_n = g(y_n)\). Poiché \(g(y) \to x_0\), allora \(x_n \to x_0\).
    Se vale l'ipotesi che \(g(y) \neq x_0\), allora \(x_n \neq x_0\), quindi possiamo applicare il limite di \(f\):
    \[f(x_{n})\to l \implies f(g(y_{n}))\to l\]
}

\ex{
    \[\lim_{ x \to 1 } \frac{\sin(x^3-1)}{x^3-1}\]
    Posto \(y=x^3 -1\), se \(x\to 1\) allora \(y\to 0\).
    \[\lim_{ y \to 0 } \frac{\sin y}{y} = 1\]
}